\section{Conclusion}
\label{sec:concl}
Using drivers' own prices, replacement intervals, resale values and daily distance from a survey of \nSurvey{} easy-bike drivers, together with NASA POWER temperature, we find that LFP is not yet the cheaper battery for the operator. At survey values a lead-acid set (BDT \costMedLA{}, \replMedLA{} months) costs BDT \tcoLA{} over ten years and an LFP set (BDT \costMedLFP{}, \replMedLFP{} months) BDT \tcoLFP{}, BDT \gapLFP{} (\gapPct\%) more, although laboratory cycle lives would suggest LFP is BDT \litGap{} cheaper. LFP matches lead-acid at a set price of about BDT \beP{} or a life of about \beL{} months. LFP still draws \gridratio\% less grid electricity, saves about \saveBike{} kg CO$_2$e per bike-year (\fleetA{} Mt per year for 1.4 M bikes) and removes \leadLow--\leadHigh{} kt of lead per year from informal recycling, at an operator support of BDT \suppFull{} billion over ten years for full conversion. Site temperature, which differs by only \Tspread\,$^\circ$C across Khulna, Chuadanga and Rajshahi, does not change the ranking, while battery heating and heat waves shorten the life of both chemistries and of lead-acid most. The answer to when lithium replaces lead is therefore economic and financial rather than climatic: it requires a cheaper or longer-lasting LFP set, or support and financing that recognise its lead and carbon benefits. The study rests on a small convenience survey, with LFP results from five drivers, and should be repeated with metered operating data and a larger, randomised sample.

\section{Reproducibility}
\label{sec:repro}
The analysis runs in MATLAB R2024a with the files in the \texttt{matlab} folder: (1) \texttt{easybike\_survey} (survey tables and figures from \texttt{survey\_cleaned.csv}); (2) \texttt{easybike\_main} (model tables and figures); (3) \texttt{easybike\_co2} (carbon, lead and policy tables); (4) \texttt{easybike\_build\_simulink} (builds \texttt{easybike\_battery.slx} and verifies it against the analytic model); (5) \texttt{easybike\_check} (compares a user's \texttt{results\_summary.csv} with \texttt{expected\_results.csv}). All inputs live in \texttt{easybike\_params.m}, each tagged with its source; the survey-based inputs are computed from \texttt{survey\_cleaned.csv}. The tables and figures in this paper were generated by a Python implementation of the same equations, which reproduced an earlier MATLAB \texttt{results\_summary.csv} to 15 significant digits. \todo{[Run easybike\_main and easybike\_check in MATLAB with the survey parameters and confirm MATCH before submission.]}
